\documentclass[10pt,a4paper]{amsart}
\usepackage[margin=19mm]{geometry}
\usepackage{amsmath,amssymb,mathtools}
\usepackage[scr=rsfs,cal=euler]{mathalfa}
\usepackage{xfrac}
\usepackage[unicode,hidelinks,bookmarks=false]{hyperref}
\newcommand{\ie}{i.e.\ }
\newcommand{\cf}{cf.\ }
\newcommand{\resp}{respectively\ }
\newcommand{\Subsection}{\subsection}
\providecommand{\nobreakdash}{\nobreak\hbox{-}}
\setlength{\emergencystretch}{2em}
\multlinegap0pt
\usepackage{bm}
\usepackage{bbm}
\usepackage{dsfont}
\usepackage{xspace}
\usepackage{cancel}
\usepackage{mathtools}
\usepackage[shortlabels]{enumitem}
\usepackage{amsthm}
\makeatletter
\@ifundefined{proposition}{\newtheorem{proposition}{Proposition}}{}
\makeatother
\title[A Bundle-Theoretic Formulation of Phonons in Crystalline Pha]{A Bundle-Theoretic Formulation of Phonons in Crystalline Phases}
\author{Aleksey Prots}
\date{Source: arXiv:2605.09323v1 (CC BY 4.0)}
\begin{document}
\maketitle

\noindent\emph{Source and licence.} A Bundle-Theoretic Formulation of Phonons in Crystalline Phases. Version arXiv:2605.09323v1 (CC BY 4.0), distributed under the Creative Commons Attribution 4.0 International licence, as recorded in the arXiv licence metadata for this exact version and shown on the abstract page https://arxiv.org/abs/2605.09323v1. Full rights notice: licenses/arxiv-2605.09323-CC-BY-4.0.txt.\par\smallskip

\noindent\textit{Editorial scope.} Counted unit: the Proposition whose source label is \texttt{prop:affine\_point\_action} in the article (source0.tex lines 255--301, source label \texttt{prop:affine\_point\_action}), delivered together with the article's own complete original proof of it (source0.tex lines 303--351, one \texttt{proof} environment of 49 lines). No mathematics is written, repaired or re-derived: statement and proof are the article's own text line for line. Required context retained uncounted: eq:cryst\_extension (lines 201--204). Every definition, display and label of those lines is retained unchanged. Omitted from this package and not redistributed: 1--200, 205--254, 302--302, 352--2185. Nothing inside a retained excerpt is omitted or altered: every retained body is delivered exactly as the article's lines are written, so no omission receipt of any kind accompanies this package. Citations: the retained excerpts carry no citation command at all, so no bibliography entry of the article is transcribed and no reference is redistributed. Reference adaptation: none was needed and none was made; every reference inside the delivered unit resolves inside it, so no unresolved reference or citation is produced. Printed slips kept literal: no printed word, symbol, hypothesis or number of the counted statement or of its proof was altered. Layout and class: the article's own class is \texttt{amsart} with packages including fontenc, inputenc, babel, lmodern, geometry, microtype, amsmath, amssymb, mathtools, amsthm, enumitem, cite, hyperref; the wrapper is the lane's pinned \texttt{amsart} editorial wrapper at 10pt on A4, which restates the theorem-like environments the retained text needs and adds the article's own macro definitions that the retained text needs (none), with extra wrapper packages (bm, bbm, dsfont, xspace, cancel, mathtools, enumitem). The delivered numbering is the unit's own. Assets: no retained span contains a figure environment, a graphics-inclusion command, a PDF-page inclusion, a table or a TikZ picture; no image file is copied into this package, the package declares no asset dependency and no image file is redistributed with it. Licence: version arXiv:2605.09323v1 of this article is distributed under the Creative Commons Attribution 4.0 International licence, as recorded in the arXiv licence metadata for this exact version and shown on the abstract page https://arxiv.org/abs/2605.09323v1. A full rights notice is delivered at licenses/arxiv-2605.09323-CC-BY-4.0.txt. Characters: every retained excerpt is pure ASCII, so the delivered engine, XeLaTeX with its default Latin Modern text font, reports no missing character.
\begin{equation}
	1\longrightarrow \Pi \longrightarrow \Gamma \longrightarrow P \longrightarrow 1.
	\label{eq:cryst_extension}
\end{equation}

\begin{proposition}[Affine action of the point group on the translation torus]
	\label{prop:affine_point_action}
	Let \(\Gamma\subset \mathbb R^d\rtimes O(d)\) be a crystallographic group. 
	Let
	\[
	\Pi=\Gamma\cap\mathbb R^d
	\]
	be its translation lattice, and let
	\[
	P=\Gamma/\Pi
	\]
	be its point group. 
	For each \(p\in P\), choose a representative
	\[
	\gamma_p=(a_p,A_p)\in\Gamma.
	\]
	Then the following statements hold.
	
	\begin{enumerate}[label=\textup{(\roman*)}]
		\item The map \(p\mapsto A_p\) defines a representation
		\[
		A:P\to \mathrm{Aut}(\Pi)\subset O(d).
		\]
		
		\item The element
		\begin{equation}
			c(p,q):=a_p+A_pa_q-a_{pq}
			\label{eq:extension_cocycle}
		\end{equation}
		belongs to \(\Pi\) and defines a \(2\)-cocycle of the group \(P\) with values in the \(P\)-module \(\Pi_A\), where the action of \(P\) on \(\Pi\) is given by the representation \(A\).
		
		\item Changing the representatives \(\gamma_p\) changes \(c\) by a coboundary. Hence the class
		\[
		[c]\in H^2(P,\Pi_A)
		\]
		is an invariant of the extension \eqref{eq:cryst_extension}. The extension splits, that is \(\Gamma\simeq\Pi\rtimes P\), if and only if this class is zero.
		
		\item The formula
		\begin{equation}
			\rho(p)[v]=[A_pv+a_p],
			\qquad [v]\in T_\Pi=\mathbb R^d/\Pi,
			\label{eq:affine_torus_action}
		\end{equation}
		defines a left affine action of \(P\) on the torus \(T_\Pi\). 
		In the symmorphic case, after a suitable choice of origin and representatives, one has \(a_p=0\), and \(\rho\) becomes the linear action \([v]\mapsto[A_pv]\).
	\end{enumerate}
\end{proposition}

\begin{proof}
	For \(\lambda\in\Pi\), conjugation by \(\gamma_p=(a_p,A_p)\) gives
	\[
	\gamma_p(\lambda,I)\gamma_p^{-1}=(A_p\lambda,I)\in\Pi.
	\]
	Hence \(A_p\Pi=\Pi\). Moreover, replacing \(\gamma_p\) by \((b_p,I)\gamma_p\), \(b_p\in\Pi\), changes only the translation part \(a_p\), not the orthogonal part \(A_p\). Thus \(p\mapsto A_p\) is well defined as a map \(P\to\mathrm{Aut}(\Pi)\subset O(d)\).
	
	For \(p,q\in P\), multiplication in \(\mathbb R^d\rtimes O(d)\) gives
	\[
	\gamma_p\gamma_q
	=
	(a_p+A_pa_q,A_pA_q).
	\]
	Since \(\gamma_p\gamma_q\) and \(\gamma_{pq}\) project to the same element \(pq\in P\), the element \(\gamma_p\gamma_q\gamma_{pq}^{-1}\) lies in \(\Pi\). Its orthogonal part is therefore the identity, so \(A_pA_q=A_{pq}\). Thus \(A:P\to\mathrm{Aut}(\Pi)\) is a representation. Its translation part gives
	\[
	c(p,q):=a_p+A_pa_q-a_{pq}\in\Pi.
	\]
	Associativity of multiplication in \(\Gamma\) gives
	\[
	A_pc(q,r)-c(pq,r)+c(p,qr)-c(p,q)=0,
	\]
	hence \(c\in Z^2(P,\Pi_A)\).
	
	If the representatives are changed to \(\gamma'_p=(b_p,I)\gamma_p\), \(b_p\in\Pi\), then \(a'_p=b_p+a_p\), and the corresponding cocycle is
	\[
	c'(p,q)
	=
	c(p,q)+b_p+A_pb_q-b_{pq}.
	\]
	Thus \(c'\) and \(c\) differ by a coboundary, so the class \([c]\in H^2(P,\Pi_A)\) is independent of the chosen representatives.
	
	The extension \eqref{eq:cryst_extension} splits if and only if \([c]=0\). Indeed, if \([c]=0\), the representatives may be changed so that \(c(p,q)=0\) for all \(p,q\). Then \(\gamma_p\gamma_q=\gamma_{pq}\), and the representatives form a subgroup of \(\Gamma\) isomorphic to \(P\). Conversely, if the extension splits, representatives can be chosen inside such a subgroup, and then \(c=0\).
	
	It remains to check the action on \(T_\Pi\). The formula
	\[
	\rho(p)[v]=[A_pv+a_p]
	\]
	is well defined with respect to \(v\): if \(v\) is replaced by \(v+\lambda\), \(\lambda\in\Pi\), then \(A_p\lambda\in\Pi\). It is also independent of the chosen representative \(\gamma_p\), since replacing \(a_p\) by \(a_p+b_p\), \(b_p\in\Pi\), does not change the class in \(\mathbb R^d/\Pi\). Finally,
	\[
	\begin{aligned}
		\rho(p)\rho(q)[v]
		&=[A_p(A_qv+a_q)+a_p]  \\
		&=[A_{pq}v+a_{pq}+c(p,q)] \\
		&=[A_{pq}v+a_{pq}]
		=\rho(pq)[v],
	\end{aligned}
	\]
	because \(c(p,q)\in\Pi\) vanishes in the quotient. Therefore \(\rho\) is a left affine action of \(P\) on \(T_\Pi\).
\end{proof}

\end{document}
